\section{Week 15: Modeling in Cognitive Science II: Reinforcement
Learning, Bayesian Cognition and AI}

\begin{quote}
\textbf{Chapter 6: Examples of modeling in different areas of cognitive
science} · Lecture 1 h · Lab 2 h \textbf{Running data:} simulated
choices in a two-armed bandit task with a reversal; a simulated
visual--haptic cue-combination experiment.
\end{quote}

\begin{center}\rule{0.5\linewidth}{0.5pt}\end{center}

\subsection{Key terms}

Important terms of this lesson, with their Persian equivalents (repeated
even if they appeared in earlier weeks).

\begin{longtable}[]{@{}
  >{\raggedright\arraybackslash}p{(\linewidth - 4\tabcolsep) * \real{0.3333}}
  >{\raggedright\arraybackslash}p{(\linewidth - 4\tabcolsep) * \real{0.3333}}
  >{\raggedright\arraybackslash}p{(\linewidth - 4\tabcolsep) * \real{0.3333}}@{}}
\caption{\centering Key terms}\tabularnewline
\toprule\noalign{}
\begin{minipage}[b]{\linewidth}\raggedright
English
\end{minipage} & \begin{minipage}[b]{\linewidth}\raggedright
\RL{\CMfa فارسی}
\end{minipage} & \begin{minipage}[b]{\linewidth}\raggedright
\RL{\CMfa توضیح کوتاه}
\end{minipage} \\
\midrule\noalign{}
\endfirsthead
\toprule\noalign{}
\begin{minipage}[b]{\linewidth}\raggedright
English
\end{minipage} & \begin{minipage}[b]{\linewidth}\raggedright
\RL{\CMfa فارسی}
\end{minipage} & \begin{minipage}[b]{\linewidth}\raggedright
\RL{\CMfa توضیح کوتاه}
\end{minipage} \\
\midrule\noalign{}
\endhead
\bottomrule\noalign{}
\endlastfoot
simulation & \RL{\CMfa شبیه‌سازی} & \RL{\CMfa تولید داده‌ی مصنوعی با اجرای مدل} \\
bandit task & \RL{\CMfa تکلیف} bandit & \RL{\CMfa انتخاب مکرر بین گزینه‌هایی با پاداش
احتمالی} \\
cue combination & \RL{\CMfa ترکیب نشانه‌ها} (cues) & \RL{\CMfa ترکیب اطلاعات چند حس
برای یک برآورد} \\
Q-learning & Q-learning & \RL{\CMfa قاعده‌ی یادگیری ارزش گزینه‌ها} \\
softmax / inverse temperature (\ensuremath{\beta}) & softmax / \RL{\CMfa دمای معکوس} &
\RL{\CMfa تبدیل ارزش‌ها به احتمال انتخاب} \\
learning rate (\ensuremath{\alpha}) & \RL{\CMfa نرخ یادگیری} & \RL{\CMfa اندازه‌ی به‌روزرسانی بعد از هر
خطا} \\
temperature & \RL{\CMfa دما} (temperature) & \RL{\CMfa پارامتر نویز در انتخاب}
token\RL{\CMfa ؛} \ensuremath{\beta} = 1/T \\
trial & \RL{\CMfa کوشش} (trial) & \RL{\CMfa یک بار انجام تکلیف توسط شرکت‌کننده} \\
likelihood & \RL{\CMfa درست‌نمایی} & \RL{\CMfa احتمال داده به‌عنوان تابعی از
پارامترها} \\
fit / fitting & \RL{\CMfa برازش} & \RL{\CMfa تنظیم پارامترها تا مدل با داده جور
شود} \\
parameter recovery & \RL{\CMfa بازیابی پارامتر} & \RL{\CMfa آیا برازش، پارامترهای
معلوم شبیه‌سازی را برمی‌گرداند} \\
parameter & \RL{\CMfa پارامتر} & \RL{\CMfa کمیت قابل تنظیم مدل که بین افراد یا شرایط
فرق می‌کند} \\
ideal observer & \RL{\CMfa مشاهده‌گر ایده‌آل} & \RL{\CMfa مدل استنباط بهینه با اطلاعات
موجود} \\
cognitive model & \RL{\CMfa مدل شناختی} & \RL{\CMfa توصیف دقیق و ریاضی یا محاسباتی
از یک فرایند ذهنی} \\
exploration / exploitation & \RL{\CMfa کاوش} / \RL{\CMfa بهره‌برداری} & \RL{\CMfa امتحان
گزینه‌های دیگر} / \RL{\CMfa انتخاب بهترین گزینه‌ی فعلی} \\
prediction error & \RL{\CMfa خطای پیش‌بینی} & \RL{\CMfa تفاوت پاداش گرفته‌شده و پاداش
مورد انتظار} \\
reinforcement learning (RL) & \RL{\CMfa یادگیری تقویتی} & \RL{\CMfa یادگیری از پاداش
و پیامد} \\
noise & \RL{\CMfa نویز} & \RL{\CMfa تغییرات تصادفی در شواهد یا در داده} \\
estimate / estimator & \RL{\CMfa برآورد} / \RL{\CMfa برآوردگر} & \RL{\CMfa مقدار به‌دست‌آمده
برای پارامتر} / \RL{\CMfa روش به دست آوردن آن} \\
reliability & \RL{\CMfa پایایی (قابلیت اطمینان)} & \RL{\CMfa دقت یک نشانه؛ عکس
واریانس آن} \\
large language model (LLM) & \RL{\CMfa مدل بزرگ زبانی} & \RL{\CMfa شبکه‌ی عصبی‌ای که
کلمه‌ی بعدی متن را پیش‌بینی می‌کند} \\
parameter trade-off & \RL{\CMfa بده‌بستان پارامترها} & \RL{\CMfa تغییر یک پارامتر که
با پارامتر دیگر جبران می‌شود} \\
RLHF / reward model & RLHF / \RL{\CMfa مدل پاداش} & \RL{\CMfa آموزش} chatbot \RL{\CMfa از
روی ترجیحات انسانی} \\
Bradley--Terry model & \RL{\CMfa مدل} Bradley--Terry & \RL{\CMfa مدل انتخاب برای
مقایسه‌های جفتی} \\
\end{longtable}

\begin{center}\rule{0.5\linewidth}{0.5pt}\end{center}

\subsection{Learning outcomes}

By the end of this week, students will be able to:

\begin{enumerate}
\def\labelenumi{\arabic{enumi}.}
\tightlist
\item
  \textbf{Formalize} reward learning with the Rescorla--Wagner /
  Q-learning rule and a softmax choice rule, and \textbf{interpret} the
  learning rate and inverse temperature.
\item
  \textbf{Write} the trial-by-trial likelihood of a learning model,
  \textbf{fit} it by MLE, and \textbf{evaluate} parameter recovery and
  model recovery against simpler strategies.
\item
  \textbf{Derive} the predictions of a Bayesian ideal observer for cue
  combination, and \textbf{test} whether simulated behavior is optimal.
\item
  \textbf{Distinguish} Bayesian models of cognition (a claim about the
  mind) from Bayesian data analysis (a method for scientists).
\item
  \textbf{Discuss} how cognitive models and AI systems inform each
  other, and how the evaluation tools of this course apply to AI systems
  proposed as models of cognition.
\end{enumerate}

\begin{center}\rule{0.5\linewidth}{0.5pt}\end{center}

\subsection{Motivation \& big picture}

\textbf{Learning from consequences.} A participant repeatedly chooses
between two slot machines, one paying off more often than the other.
Halfway through, the machines secretly swap. How quickly does the
participant notice and adapt? How much do they explore versus exploit?
One of the most successful models in all of psychology and neuroscience
answers with a single equation: update your expectation by a fraction of
the \textbf{prediction error}, the difference between what you got and
what you expected. The same error-driven rule appears in classical
conditioning (Rescorla--Wagner), in connectionist learning (the delta
rule), in the firing of dopamine neurons, and in the
reinforcement-learning algorithms behind modern AI systems. Fitting it
to choice data is one of the most common uses of cognitive modeling
today, in computational psychiatry for instance.

\textbf{The mind as a statistician.} When you see and touch an object at
the same time, vision and touch both give noisy estimates of its size. A
\emph{Bayesian ideal observer} combines them by weighting each cue by
its reliability, which gives an estimate more precise than either alone.
Do people do this? This is a different use of the word ``Bayesian'' from
Weeks 9--13: here Bayes' rule is a \textbf{hypothesis about the mind},
at Marr's computational level (Week 1).

\textbf{Cognitive models \ensuremath{\leftrightarrow} AI.} Many ideas travel in both directions
between cognitive modeling and AI: error-driven learning, neural
networks, reinforcement learning, Bayesian inference. Today AI systems,
including large language models, are themselves proposed as models of
human cognition. The tools of this course (likelihoods, model
comparison, generalization tests, recovery) are exactly what is needed
to evaluate such claims.

\paragraph{Where is this used?}

\begin{itemize}
\tightlist
\item
  \textbf{Computational psychiatry}: reinforcement-learning models are
  fitted to choices of patients with depression, addiction,
  schizophrenia or OCD, to test whether they learn differently from
  rewards and punishments (e.g.~a lower learning rate after losses).
\item
  \textbf{Neuroscience}: the prediction error of the delta rule matches
  the activity of dopamine neurons, one of the best-known links between
  a cognitive model and the brain.
\item
  \textbf{Education and games}: adaptive learning systems and game
  design use RL ideas to model how players and learners adjust to
  feedback.
\item
  \textbf{Perception and robotics}: optimal cue combination explains how
  people combine vision, touch and hearing, and the same weighting rule
  is used in sensor fusion.
\item
  \textbf{AI}: modern reinforcement learning and the training of
  chatbots from human preferences (RLHF) use the same choice models (see
  the AI \& LLMs section).
\end{itemize}

\begin{center}\rule{0.5\linewidth}{0.5pt}\end{center}

\Needspace{11\baselineskip}

\subsection{Lecture notes}

\paragraph{In plain words}

\textbf{Learning from rewards.} Each option (e.g.~each of two slot
machines) has a ``value'' in the participant's head: how much reward
they expect from it. After choosing an option and seeing the outcome,
the value is nudged toward what actually happened: \emph{new value = old
value + learning rate × (reward − old value)}. The term in brackets is
the \textbf{prediction error}, the surprise. The \textbf{learning rate}
\(\alpha\) says how big the nudge is. Choices are not perfectly greedy:
the \textbf{softmax} rule picks the better-valued option more often, and
the \textbf{inverse temperature} \(\beta\) sets how strongly (low
\(\beta\) = more random exploring, high \(\beta\) = almost always the
best). Fitting this model means finding the \(\alpha\) and \(\beta\)
that make the participant's actual sequence of choices most likely.
\textbf{The Bayesian mind:} when two senses give noisy estimates of the
same thing (seeing and touching an object's size), the best guess is a
weighted average that trusts the more reliable sense more. People often
come close to this optimal weighting; the model predicts it with no free
parameters, which makes it a strong test.

\textbf{Suggested timing (60 min):} reward learning 20 · fitting
learning models 10 · Bayesian models of cognition 15 · cognitive models
and AI 10 · course synthesis 5.

\subsubsection{1. Reinforcement learning: the delta rule and Q-learning}

\textbf{Rescorla--Wagner / delta rule.} After choosing option \(o\) on
trial \(t\) and receiving reward \(r_t \in \{0, 1\}\):

\[
\delta_t = r_t - Q_t(o) \quad\text{(prediction error)}, \qquad
Q_{t+1}(o) = Q_t(o) + \alpha\,\delta_t,
\]

and the values of unchosen options stay the same. The \textbf{learning
rate} \(\alpha \in [0, 1]\) sets how strongly new outcomes overwrite old
expectations. A high \(\alpha\) adapts quickly but is noisy; a low
\(\alpha\) is stable but slow to notice a reversal.

\textbf{Softmax choice rule.} Values are turned into choice
probabilities with \textbf{inverse temperature} \(\beta \ge 0\):

\[
P(\text{choose } o \text{ on trial } t) = \frac{\exp(\beta\, Q_t(o))}{\sum_{o'} \exp(\beta\, Q_t(o'))}.
\]

\(\beta = 0\) means random choice (pure exploration); a large \(\beta\)
means always picking the currently best option (exploitation). For two
options this is a logistic function of the value difference:
\(P(o_1) = 1/(1 + e^{-\beta(Q_1 - Q_2)})\).

\textbf{Connections:} - The update is the \textbf{delta rule} of
connectionist networks (Verguts, 2022): a single-layer network trained
by error correction. - Prediction errors \(\delta_t\) resemble the
phasic responses of midbrain dopamine neurons, which made RL models
central to computational neuroscience. - Q-learning with
temporal-difference errors is the basis of modern reinforcement learning
in AI.

\paragraph{Reading the math}

\begin{itemize}
\tightlist
\item
  \(Q_t(o)\): the value of option \(o\) on trial \(t\) (the subscript
  \(t\) is the trial number).
\item
  \(\delta_t = r_t - Q_t(o)\): reward received minus reward expected.
  Positive = pleasant surprise, negative = disappointment.
\item
  \(Q_{t+1}(o) = Q_t(o) + \alpha\,\delta_t\). \textbf{Example:}
  \(Q = 0.5\), \(\alpha = 0.3\), reward 1: \(\delta = 0.5\), new
  \(Q = 0.5 + 0.15 = 0.65\). With no reward: \(\delta = -0.5\), new
  \(Q = 0.35\).
\item
  \textbf{Softmax for two options:}
  \(P(\text{option 1}) = 1/(1 + e^{-\beta(Q_1 - Q_2)})\). With
  \(Q_1 - Q_2 = 0.25\): \(\beta = 0\) gives 0.50 (random), \(\beta = 3\)
  gives 0.68, \(\beta = 20\) gives 0.99 (almost always option 1).
\end{itemize}

\subsubsection{2. Fitting a learning model}

\textbf{The likelihood is sequential.} Values depend on the
participant's own past choices and rewards, so we run the model through
the participant's actual trial sequence and multiply the probabilities
it assigned to the choices actually made:

\[
\ell(\alpha, \beta) = \sum_{t=1}^{T} \ln P\big(\text{choice}_t \mid Q_t(\cdot;\ \alpha, \text{history}_{<t}),\ \beta\big).
\]

\textbf{Practical points:} - Reparameterize:
\(\alpha = 1/(1 + e^{-\eta_\alpha})\) and \(\beta = e^{\eta_\beta}\)
(Week 3), and use multi-start. - \textbf{\(\alpha\)--\(\beta\)
trade-off:} noisy choice (low \(\beta\)) and slow or erratic learning
(extreme \(\alpha\)) can produce similar-looking behavior, so the
estimates are often correlated across fits, and the sign of the
correlation depends on the design and the parameter range. A further
problem is that very large \(\beta\) values are hard to estimate at all
(lab Part C). Recovery studies at the actual design are essential (Week
7). - \textbf{The design matters:} reversals and enough trials are
needed to estimate \(\alpha\); stable reward probabilities give little
information about it. - \textbf{Alternative strategies} must be
compared: e.g., \emph{win-stay/lose-shift} (WSLS: repeat a rewarded
choice, switch after no reward, with lapse rate \(\varepsilon\)), or a
biased-random model. A good RL fit means little if WSLS fits just as
well.

\paragraph{Reading the math}

\(\ell(\alpha, \beta) = \sum_t \ln P(\text{choice}_t \mid Q_t, \beta)\):
run the learning model through the participant's \emph{own} sequence of
choices and rewards, trial by trial. On each trial, compute the
probability the model gave to the choice the participant actually made,
take its log, and add up. The worked example (Part 1) does this for four
trials: \(\ln 0.5 + \ln 0.679 + \ln 0.593 + \ln 0.755 = -1.884\). The
best \((\alpha, \beta)\) maximize this sum. It is Week 4's likelihood,
except that each trial's probability depends on the history before it.

\subsubsection{3. Bayesian models of cognition}

\textbf{Cue combination.} A visual cue
\(x_V \sim \mathcal{N}(S, \sigma_V)\) and a haptic cue
\(x_H \sim \mathcal{N}(S, \sigma_H)\) both measure the true size \(S\).
With a flat prior, the posterior over \(S\) is normal with mean

\[
\hat S = w_V x_V + w_H x_H, \qquad w_V = \frac{1/\sigma_V^2}{1/\sigma_V^2 + 1/\sigma_H^2}, \quad w_H = 1 - w_V,
\]

and SD

\[
\sigma_{VH} = \sqrt{\frac{\sigma_V^2\,\sigma_H^2}{\sigma_V^2 + \sigma_H^2}} \;\le\; \min(\sigma_V, \sigma_H).
\]

\textbf{Testable predictions} (tested, e.g., by Ernst \& Banks, 2002):
(1) the weights should track the cues' reliabilities, which the
experimenter can manipulate by degrading one cue; (2) bimodal precision
should be \emph{better} than the best single cue, by exactly the amount
above. The single-cue SDs are measured separately, so the ideal-observer
predictions have \textbf{no free parameters}. That makes them strong,
falsifiable tests.

\textbf{What a Bayesian model of cognition claims.} That behavior
approximates optimal inference given the agent's noise, prior and goals:
a \textbf{computational-level} account (Marr). It is silent about the
algorithm: the brain might approximate Bayes with sampling, heuristics
or neural population codes.

\textbf{Criticisms to keep in mind:} with free choices of priors,
likelihoods and loss functions, almost any behavior can be called
``optimal'' after the fact (the flexibility problem of Week 6). Strong
tests fix these components independently, as in cue combination.

\textbf{Two meanings, one more time:}

\begin{longtable}[]{@{}
  >{\raggedright\arraybackslash}p{(\linewidth - 4\tabcolsep) * \real{0.3333}}
  >{\raggedright\arraybackslash}p{(\linewidth - 4\tabcolsep) * \real{0.3333}}
  >{\raggedright\arraybackslash}p{(\linewidth - 4\tabcolsep) * \real{0.3333}}@{}}
\caption{\centering 3. Bayesian models of cognition}\tabularnewline
\toprule\noalign{}
\begin{minipage}[b]{\linewidth}\raggedright
\end{minipage} & \begin{minipage}[b]{\linewidth}\raggedright
Bayesian model of cognition
\end{minipage} & \begin{minipage}[b]{\linewidth}\raggedright
Bayesian data analysis
\end{minipage} \\
\midrule\noalign{}
\endfirsthead
\toprule\noalign{}
\begin{minipage}[b]{\linewidth}\raggedright
\end{minipage} & \begin{minipage}[b]{\linewidth}\raggedright
Bayesian model of cognition
\end{minipage} & \begin{minipage}[b]{\linewidth}\raggedright
Bayesian data analysis
\end{minipage} \\
\midrule\noalign{}
\endhead
\bottomrule\noalign{}
\endlastfoot
Who does the inference? & The participant's mind (hypothetically) & The
scientist \\
Prior represents & The participant's beliefs about the world & The
scientist's knowledge about parameters \\
Can be wrong if\ldots{} & People are not optimal & The model or priors
are poorly chosen \\
Examples in this course & Cue combination (today) & Weeks 9--13 \\
\end{longtable}

You can analyze a Bayesian model of cognition with maximum likelihood,
and a non-Bayesian model (e.g., Q-learning) with Bayesian data analysis.

\paragraph{Reading the math}

\begin{itemize}
\tightlist
\item
  \(w_V = \dfrac{1/\sigma_V^2}{1/\sigma_V^2 + 1/\sigma_H^2}\): the
  weight of vision is its precision (1/variance) divided by the total
  precision. A sense twice as precise (half the SD, a quarter of the
  variance) gets four times the weight.
\item
  \(\sigma_{VH} = \sqrt{\sigma_V^2\sigma_H^2/(\sigma_V^2 + \sigma_H^2)}\):
  the combined SD, always smaller than the better single sense.
  \textbf{Example:} \(\sigma_V = 1\), \(\sigma_H = 2\):
  \(w_V = 1/(1 + 1/4) = 0.8\) and \(\sigma_{VH} = \sqrt{4/5} = 0.894\),
  about 11\% better than vision alone.
\item
  \textbf{The test:} measure \(\sigma_V\) and \(\sigma_H\) separately
  (one sense at a time), predict \(w_V\) and \(\sigma_{VH}\) with no
  free parameters, then check the two-sense data.
\end{itemize}

\paragraph{For an engineer}

\begin{itemize}
\tightlist
\item
  \textbf{The delta rule is the LMS update.}
  \(Q \leftarrow Q + \alpha(r - Q)\) is a one-tap LMS (or an
  exponentially weighted moving average with forgetting factor
  \(1-\alpha\)): large \(\alpha\) tracks changes quickly but is noisy,
  small \(\alpha\) is smooth but slow, the familiar
  tracking--misadjustment trade-off.
\item
  \textbf{Softmax} is the Boltzmann/Gibbs distribution; \(\beta\) is an
  inverse temperature in the physical sense. Sampling with a temperature
  is also how language models generate text (Week 1 AI box).
\item
  \textbf{Fitting \(\alpha\) and \(\beta\) from choices} is system
  identification of an adaptive system from its binary outputs only: you
  never observe \(Q\), only the decisions it produces. That is why
  \(\alpha\) and \(\beta\) trade off and why recovery checks are
  essential.
\item
  \textbf{Optimal cue combination} uses exactly the inverse-variance
  weights of \textbf{maximal-ratio combining} and of sensor fusion in a
  Kalman filter: the combined estimate has variance
  \((1/\sigma_V^2 + 1/\sigma_H^2)^{-1}\). The scientific question is
  whether the \emph{brain} implements this fusion.
\end{itemize}

\subsubsection{4. Cognitive models and artificial intelligence}

\begin{longtable}[]{@{}
  >{\raggedright\arraybackslash}p{(\linewidth - 2\tabcolsep) * \real{0.5000}}
  >{\raggedright\arraybackslash}p{(\linewidth - 2\tabcolsep) * \real{0.5000}}@{}}
\caption{\centering 4. Cognitive models and artificial intelligence}\tabularnewline
\toprule\noalign{}
\begin{minipage}[b]{\linewidth}\raggedright
Direction
\end{minipage} & \begin{minipage}[b]{\linewidth}\raggedright
Examples
\end{minipage} \\
\midrule\noalign{}
\endfirsthead
\toprule\noalign{}
\begin{minipage}[b]{\linewidth}\raggedright
Direction
\end{minipage} & \begin{minipage}[b]{\linewidth}\raggedright
Examples
\end{minipage} \\
\midrule\noalign{}
\endhead
\bottomrule\noalign{}
\endlastfoot
Cognitive science \ensuremath{\rightarrow} AI & Neural networks (inspired by neurons and
learning rules); reinforcement learning (from animal learning theory);
exemplar/prototype ideas in machine-learning classifiers
(nearest-neighbor vs.~centroid methods) \\
AI \ensuremath{\rightarrow} cognitive science & Deep networks as models of visual cortex; RL
agents as models of decision making; large language models proposed as
models of language processing and of human-like reasoning \\
Shared mathematics & Likelihoods and loss functions; gradient-based
optimization (Week 3); Bayesian inference and MCMC; cross-validation and
generalization (Week 6) \\
\end{longtable}

\textbf{Evaluating an AI system as a cognitive model} uses exactly the
course toolkit: - Does it predict human data \emph{beyond} fitting it,
e.g., new conditions and individual differences (generalization, Week
6)? - Does it beat simpler, interpretable cognitive models on a fair
comparison (Weeks 7, 10)? - Can we tell from data whether humans and the
model use the same mechanism, or only produce similar outputs (model
mimicry and recovery, Week 7)? - Are its parameters or internal states
interpretable as psychological constructs (Week 1)?

A model with billions of parameters that fits behavior well is not
automatically an \emph{explanation}: flexibility, identifiability and
interpretability matter as much as fit.

\subsubsection{5. Course synthesis}

Across very different domains (forgetting, detection, categorization,
decisions, learning, perception), the workflow has been the same:
\textbf{formalize \ensuremath{\rightarrow} simulate \ensuremath{\rightarrow} fit \ensuremath{\rightarrow} evaluate \ensuremath{\rightarrow} recover \ensuremath{\rightarrow} interpret.}
The domain changes the model; the logic of modeling stays.

\paragraph{For the instructor}

\begin{itemize}
\tightlist
\item
  \textbf{Play the bandit game with the class.} Two ``machines'' (two
  students with hidden coins or a random number app, one paying 80\% of
  the time and one 20\%), and halfway through secretly swap them. The
  class experiences exploration, exploitation and reversal before seeing
  any equation.
\item
  \textbf{Do the four Q-learning trials of the worked example on the
  board} with a table: value before, choice, reward, prediction error,
  value after.
\item
  \textbf{For cue combination, use your engineering intuition openly}:
  ``combining two noisy sensors'' is a natural way to present it, and
  students from psychology find it illuminating.
\item
  \textbf{Keep the AI discussion grounded.} Section 4 can trigger long
  debates; use the four evaluation questions (generalization, fair
  comparison, mimicry, interpretability) to keep them focused.
\item
  \textbf{Typical questions.}

  \begin{itemize}
  \tightlist
  \item
    \emph{``Why is \(\beta\) hard to estimate?''} Once choices are
    almost deterministic, increasing \(\beta\) further hardly changes
    the likelihood.
  \item
    \emph{``Is the brain really doing Q-learning?''} It is a good
    description of many choice data and matches dopamine signals, but
    more complex models (model-based learning) also have support.
  \item
    \emph{``Do people always combine cues optimally?''} Often close, but
    not always; the lab participant over-weights touch.
  \end{itemize}
\end{itemize}

\begin{center}\rule{0.5\linewidth}{0.5pt}\end{center}

\Needspace{11\baselineskip}

\subsection{Worked example}

\subsubsection{Part 1: four trials of Q-learning by hand}

Two options; \(\alpha = 0.5\), \(\beta = 3\); initial values
\(Q(1) = Q(2) = 0.5\).

\begin{landscape}\small
\begin{longtable}[]{@{}
  >{\raggedright\arraybackslash}p{(\linewidth - 14\tabcolsep) * \real{0.1250}}
  >{\raggedright\arraybackslash}p{(\linewidth - 14\tabcolsep) * \real{0.1250}}
  >{\raggedright\arraybackslash}p{(\linewidth - 14\tabcolsep) * \real{0.1250}}
  >{\raggedright\arraybackslash}p{(\linewidth - 14\tabcolsep) * \real{0.1250}}
  >{\raggedright\arraybackslash}p{(\linewidth - 14\tabcolsep) * \real{0.1250}}
  >{\raggedright\arraybackslash}p{(\linewidth - 14\tabcolsep) * \real{0.1250}}
  >{\raggedright\arraybackslash}p{(\linewidth - 14\tabcolsep) * \real{0.1250}}
  >{\raggedright\arraybackslash}p{(\linewidth - 14\tabcolsep) * \real{0.1250}}@{}}
\caption{\centering Part 1: four trials of Q-learning by hand}\tabularnewline
\toprule\noalign{}
\begin{minipage}[b]{\linewidth}\raggedright
Trial
\end{minipage} & \begin{minipage}[b]{\linewidth}\raggedright
\(Q(1)\)
\end{minipage} & \begin{minipage}[b]{\linewidth}\raggedright
\(Q(2)\)
\end{minipage} & \begin{minipage}[b]{\linewidth}\raggedright
\(P(\text{choose } 1)\)
\end{minipage} & \begin{minipage}[b]{\linewidth}\raggedright
Choice
\end{minipage} & \begin{minipage}[b]{\linewidth}\raggedright
Reward
\end{minipage} & \begin{minipage}[b]{\linewidth}\raggedright
\(\delta\)
\end{minipage} & \begin{minipage}[b]{\linewidth}\raggedright
Update
\end{minipage} \\
\midrule\noalign{}
\endfirsthead
\toprule\noalign{}
\begin{minipage}[b]{\linewidth}\raggedright
Trial
\end{minipage} & \begin{minipage}[b]{\linewidth}\raggedright
\(Q(1)\)
\end{minipage} & \begin{minipage}[b]{\linewidth}\raggedright
\(Q(2)\)
\end{minipage} & \begin{minipage}[b]{\linewidth}\raggedright
\(P(\text{choose } 1)\)
\end{minipage} & \begin{minipage}[b]{\linewidth}\raggedright
Choice
\end{minipage} & \begin{minipage}[b]{\linewidth}\raggedright
Reward
\end{minipage} & \begin{minipage}[b]{\linewidth}\raggedright
\(\delta\)
\end{minipage} & \begin{minipage}[b]{\linewidth}\raggedright
Update
\end{minipage} \\
\midrule\noalign{}
\endhead
\bottomrule\noalign{}
\endlastfoot
1 & 0.50 & 0.50 & 0.500 & 1 & 1 & \(1 - 0.5 = +0.5\) &
\(Q(1) = 0.5 + 0.5 \times 0.5 = 0.75\) \\
2 & 0.75 & 0.50 & \(1/(1 + e^{-3 \times 0.25}) = 0.679\) & 1 & 0 &
\(0 - 0.75 = -0.75\) & \(Q(1) = 0.75 - 0.375 = 0.375\) \\
3 & 0.375 & 0.50 & \(1/(1 + e^{-3 \times (-0.125)}) = 0.407\) & 2 & 1 &
\(1 - 0.5 = +0.5\) & \(Q(2) = 0.75\) \\
4 & 0.375 & 0.75 & \(1/(1 + e^{-3 \times (-0.375)}) = 0.245\) & 2 & 1 &
\(+0.25\) & \(Q(2) = 0.875\) \\
\end{longtable}
\end{landscape}

\textbf{Log-likelihood of these four choices:}
\(\ln 0.5 + \ln 0.679 + \ln(1 - 0.407) + \ln(1 - 0.245) = -0.693 - 0.387 - 0.523 - 0.281 = -1.884\).

\textbf{Interpretation.} With \(\alpha = 0.5\) a single non-reward
halves the value of option 1, and the model predicts a switch, which the
participant made. The likelihood rewards parameter values that predict
the participant's actual choice sequence.

\subsubsection{Part 2: an ideal observer}

Visual SD \(\sigma_V = 1\) mm, haptic SD \(\sigma_H = 2\) mm. Weights:
\(w_V = 1/(1 + 1/4) = 0.8\), \(w_H = 0.2\). Predicted combined SD:
\(\sqrt{(1 \times 4)/(1 + 4)} = 0.894\) mm, about 11\% better than
vision alone. If \(x_V = 50\) and \(x_H = 54\) mm, the optimal estimate
is \(0.8 \times 50 + 0.2 \times 54 = 50.8\) mm.

\textbf{Interpretation.} An observer who simply \emph{averaged} the cues
(\(w = 0.5\)) would have SD
\(\sqrt{0.25 \times 1 + 0.25 \times 4} = 1.118\) mm, \emph{worse} than
vision alone. Optimal weighting is what makes combining cues worthwhile.

\begin{center}\rule{0.5\linewidth}{0.5pt}\end{center}

\Needspace{11\baselineskip}

\subsection{Lab (2 hours)}

\subsubsection{Goals}

\begin{enumerate}
\def\labelenumi{\arabic{enumi}.}
\tightlist
\item
  Simulate a Q-learning agent in a reversal bandit task.
\item
  Fit \((\alpha, \beta)\) by MLE; run parameter recovery and diagnose
  the \(\alpha\)--\(\beta\) trade-off.
\item
  Compare Q-learning with win-stay/lose-shift and biased-random models.
\item
  Test whether a simulated observer combines cues optimally.
\end{enumerate}

\textbf{Timing:} Part A 20 min · Part B 20 min · Part C 25 min · Part D
30 min · Part E 25 min.

\subsubsection{Part C: parameter recovery and the \ensuremath{\alpha}--\ensuremath{\beta} trade-off (25
min)}

\textbf{Interpret:} with 200 trials and one reversal, \(\alpha\) is
recovered well (r = 0.90) but \(\beta\) hardly at all (r = 0.17). This
also explains the Part B fit (\(\hat\beta = 3.8\) for a true 5.0). Look
at the large-\(\beta\) region of the \(\beta\) panel. Why are high
inverse temperatures hard to pin down? (Hint: once choices are almost
deterministic, increasing \(\beta\) further barely changes the
likelihood.) The recovered \(\alpha\) and \(\beta\) are also correlated
(0.31) although the true values were independent: a trade-off. What does
this imply for a study that correlates \(\beta\) with, say, a clinical
questionnaire?

\subsubsection{Part E: an ideal observer for cue combination (25 min)}

A simulated participant judges the size of objects seen and felt
simultaneously. We know the single-cue noise levels from unimodal
blocks. Does the participant combine the cues optimally?

\textbf{Interpret:} is this participant optimal? The empirical visual
weight (0.56) is well below the optimal 0.80: the participant
over-weights the less reliable haptic cue. As a result, their bimodal SD
(1.006 mm) is no better than vision alone (1.000 mm), instead of the
0.894 mm an ideal observer would achieve. Note that the ideal-observer
prediction had \textbf{no free parameters}: that is what makes this a
strong test.

(We simulated this participant with a true weight of 0.65. The estimate
0.56 is off by about 3 standard errors, an unusually unlucky sample but
a useful reminder that empirical weights also have sampling error.
Exercise 4 puts bootstrap intervals around them.)

\subsubsection{Exercises (increasing difficulty)}

Do you get −1.884 for the log-likelihood?

\textbf{Exercise 2 \ensuremath{\star}.} Rerun Part C with 50 trials (no reversal) and
with 400 trials (three reversals). How does recovery of \(\alpha\)
change? What does this imply for experimental design?

\textbf{Exercise 3 \ensuremath{\star}\ensuremath{\star} (two learning rates).} Many studies find different
learning rates for positive and negative prediction errors. Formulate and evaluate a
model with \(\alpha^{+}\) and \(\alpha^{-}\), simulate an agent with
\(\alpha^{+} = 0.4\), \(\alpha^{-} = 0.1\), and test with AIC/BIC and a
recovery study whether the extra parameter is identifiable at the Part A
design.

\textbf{Exercise 4 \ensuremath{\star}\ensuremath{\star} (optimal weights).} In Part E, simulate an
observer with \(w_V = 0.8\) (optimal) and one with \(w_V = 0.5\)
(averaging), 100 bimodal trials each. Use a bootstrap (Week 5) to put
confidence intervals around the empirical weights and SDs. With 100
trials, can you distinguish the two observers?

\textbf{Exercise 5 \ensuremath{\star}\ensuremath{\star}\ensuremath{\star} (Advanced: hierarchical RL).} Simulate 20 agents
with individual \((\alpha_j, \beta_j)\) and only 80 trials each.

\subsubsection{Lab check-off}

Show: (1) the MLE fit and the recovery plots with your interpretation of
the \(\beta\) panel, (2) the model-recovery table, and (3) your answer
on the optimality of the cue-combination observer.

\begin{center}\rule{0.5\linewidth}{0.5pt}\end{center}

\subsection{Connection to AI \& LLMs}

\textbf{The choice rule behind chatbot training.} Modern chatbots are
fine-tuned with \emph{reinforcement learning from human feedback}
(RLHF). People see two responses and say which they prefer, a
\textbf{reward model} is fitted to these choices, and the chatbot is
then trained by reinforcement learning to maximize the learned reward
(Christiano et al., 2017; Ouyang et al., 2022). The reward model is
usually the \textbf{Bradley--Terry} model,
\(P(i \succ j) = 1/(1 + e^{-(r_i - r_j)})\). This is the softmax choice
rule of this week with \(\beta = 1\), and the Luce choice rule of the
GCM. Fitting it is an MLE problem from Week 4. So human choice models
sit at the heart of how today's AI systems are aligned with human
preferences, and so do their problems: noisy choices, identifiability
(only \emph{differences} in \(r\) matter, so one value is fixed), and
the gap between what people choose and what is actually good.

With about 300 noisy comparisons, the estimated order of responses 1 and
2 is even reversed. A reward model is a fitted cognitive model, and it
needs uncertainty estimates and recovery checks like any other.

Is the reversal of responses 1 and 2 within the uncertainty? - \ensuremath{\star}\ensuremath{\star} Run a
parameter recovery: how many comparisons are needed before the order of
all four responses is recovered correctly in 90\% of simulated datasets?
- \ensuremath{\star}\ensuremath{\star} Simulate raters who prefer \emph{longer} answers regardless of
quality (add a length bonus to the true reward). What does the fitted
reward model learn, and what will the chatbot then be trained to do?

\emph{Optional reading:} Christiano, P. F., et al.~(2017). Deep
reinforcement learning from human preferences. \emph{NeurIPS}. Ouyang,
L., et al.~(2022). Training language models to follow instructions with
human feedback. \emph{NeurIPS}.

\begin{center}\rule{0.5\linewidth}{0.5pt}\end{center}

\subsection{Common mistakes \& misconceptions}

\begin{enumerate}
\def\labelenumi{\arabic{enumi}.}
\tightlist
\item
  \textbf{Fitting a learning model to aggregated choice proportions.}
  The likelihood must follow each participant's own trial sequence,
  because values depend on their own history.
\item
  \textbf{Interpreting \(\alpha\) without a recovery check.} \(\alpha\)
  and \(\beta\) trade off, and \(\alpha\) is poorly identified without
  reversals or enough trials.
\item
  \textbf{Not comparing with simple strategies.} Win-stay/lose-shift or
  perseveration can mimic learning. Show that the RL model wins a fair
  comparison.
\item
  \textbf{Forgetting that the softmax uses value \emph{differences}.}
  Adding a constant to all Q-values changes nothing, while the scale of
  the values and \(\beta\) trade off.
\item
  \textbf{``People are Bayesian'' from a post-hoc fit.} If the prior and
  likelihood were chosen to fit the behavior, optimality is assumed, not
  tested. Strong tests fix them independently (cue combination).
\item
  \textbf{Confusing the two Bayesianisms} (model of the mind vs.~method
  of analysis).
\item
  \textbf{``A large AI model that fits human data is a good cognitive
  model.''} Fit is not explanation. Ask about generalization, mimicry,
  identifiability and interpretability, as for any model.
\end{enumerate}

\begin{center}\rule{0.5\linewidth}{0.5pt}\end{center}

\subsection{Check your understanding}

\textbf{1.} Write the delta-rule update and explain the role of the
prediction error. \textbf{Answer} \(Q \leftarrow Q + \alpha(r - Q)\).
The prediction error \(r - Q\) is the difference between outcome and
expectation: positive errors raise the value, negative errors lower it,
and the learning rate scales the change.

\textbf{2.} What does the inverse temperature \(\beta\) control? What
behavior do \(\beta = 0\) and \(\beta \to \infty\) produce?
\textbf{Answer} The sensitivity of choices to value differences
(exploration vs.~exploitation). \(\beta = 0\): random choice regardless
of values. \(\beta \to \infty\): always choose the option with the
highest current value.

\textbf{3.} Why must a learning model's likelihood be computed trial by
trial on each participant's actual sequence? \textbf{Answer} The model's
predictions on trial \(t\) depend on that participant's own previous
choices and rewards (through the Q-values), so the probability of each
observed choice can only be computed by running the model through the
observed history.

\textbf{4.} Explain the \(\alpha\)--\(\beta\) trade-off. \textbf{Answer}
Different combinations of learning rate and choice noise can produce
similar choice sequences. For example, sluggish adaptation can come from
slow learning (low \(\alpha\)) or from noisy, value-insensitive choice
(low \(\beta\)). The estimates are therefore correlated and individually
uncertain; the direction of the correlation depends on the design. In
addition, high \(\beta\) values are weakly identified because
near-deterministic choices look the same for many large \(\beta\).

\textbf{5.} Vision has SD 2 and touch SD 2. What are the optimal weights
and the predicted combined SD? \textbf{Answer} Equal reliabilities give
\(w_V = w_H = 0.5\), and the combined SD is
\(\sqrt{4 \times 4/8} = \sqrt{2} = 1.41\), better than either cue alone
(2).

\textbf{6.} Give one example each of an idea that moved from cognitive
science to AI and from AI to cognitive science. \textbf{Answer}
Cognitive science \ensuremath{\rightarrow} AI: error-driven learning (Rescorla--Wagner / delta
rule) leading to neural-network training and reinforcement learning. AI
\ensuremath{\rightarrow} cognitive science: deep convolutional networks used as models of the
ventral visual stream, or RL agents and language models proposed as
models of human behavior.

\textbf{7.} A researcher shows that a large language model's responses
correlate highly with human responses in a reasoning task. What further
evidence would you ask for before accepting it as a cognitive model?
\textbf{Answer} Prediction of new conditions and of individual
differences (generalization); a fair comparison against simpler
cognitive models; tests that distinguish the model's mechanism from
alternatives that produce the same outputs (mimicry); and some
correspondence between its internal states and psychological constructs.

\begin{center}\rule{0.5\linewidth}{0.5pt}\end{center}

\subsection{Assignment}

Use this week's lab time for project work and the course wrap-up.

\begin{center}\rule{0.5\linewidth}{0.5pt}\end{center}

\subsection{Further reading}

\begin{itemize}
\tightlist
\item
  \textbf{Verguts (2022)}: the chapters on the delta rule and on
  reinforcement learning (error-driven learning and its neural basis)
  {[}verify chapters{]}. \emph{(Main reading for Part 1.)}
\item
  \textbf{Farrell \& Lewandowsky (2018)}: the chapters on neural network
  models and on models in neuroscience {[}verify chapter numbers{]}.
\item
  \textbf{Anderson (2014)}: the chapters on neural networks and
  learning, which connect cognitive and computational neuroscience
  models {[}verify chapters{]}.
\item
  \emph{Optional, outside the course list:} Ernst, M. O., \& Banks, M.
  S. (2002). Humans integrate visual and haptic information in a
  statistically optimal fashion. \emph{Nature, 415}, 429--433. Wilson,
  R. C., \& Collins, A. G. E. (2019). Ten simple rules for the
  computational modeling of behavioral data. \emph{eLife, 8}, e49547
  (its running example is a bandit task like today's).
\end{itemize}

\emph{Preview for Week 16:} final project presentations, and a closing
session on good modeling practice.
